\scanpage{77}
\mspage{\unsure{70|78}}
% archive sheet number: 76
Mais supposons maintenant que $u$ soit admissible. Alors je vais définir la flèche en sens initial (p. 68), \deleted{de la façon suivante}: Je choisis une flèche
\[
\deleted{(\alpha)\ \operatorname{Hom}_W(u^*F\unsure{,|)}}
\]
\[
i'':u^*F\longrightarrow F'
\]
avec $i''\in\mathrm{TC}\subset C\cap W$, $F'\in\operatorname{Ob}M(I')_f$, (donc en fait \deleted{\unsure{de ce fait}} $F'\in\operatorname{Ob}M(I')_{cf}$ car $u^*F$ est cofibrant, \unsure{comme on sait|comme quoi})

\deleted{ou ce qui revient au $=$ \unsure{par}} et par adjonction la flèche $i'$
\[
i':F\unsure{-|\longrightarrow}u_*F'.
\]
Je considère le diagramme (comm. !)
\diagram{
\begin{tikzcd}[ampersand replacement=\&,column sep=huge,row sep=large]
\operatorname{Hom}(F',\unsure{G'|C'})\arrow[r,"\alpha"]\arrow[d,"\lambda:"',"f'\mapsto f'\circ i''"]\&\operatorname{Hom}(u_*F',u_*G')\arrow[d,"\mu:f\mapsto f\circ i'"']\\
\operatorname{Hom}(u^*F,G')\arrow[r,"\mathrm{adj.}"']\&\operatorname{Hom}(F,u_*G')
\end{tikzcd}
}
On voit que \deleted{$\unsure{F'}$} l'on a encore $F'$ cofibrant (car $u^*F$ l'est et $\unsure{i'|i''}\in\mathrm{TC}\subset C$) donc $\operatorname{Hom}_W(F',G')$ est un quotient de $\operatorname{Hom}(F',G')$, et $\lambda$ induit un isomorphisme
\[
\operatorname{Hom}_W(F',G')\xrightarrow{\sim}\operatorname{Hom}_W(u^*F,G')
\]
(car $i''\in\mathrm{TC}\deleted{\unsure{)|f}}\subset W$). Comme $F',G'\in\operatorname{Ob}M(I')_f$, et que $u_*^f$ est compatible avec $W(I')_f$, $W(I)_f$, il \unsure{s'ensuit} que la flèche horizontale $\alpha$
